\section{Related Work and Current Methods}
\label{sec:related}

\paragraph{Domain adaptation in clinical and general time series.}
Domain adaptation for electronic health records has been approached through
representation-space alignment~\cite{purushotham2017vrada,zhang2022adadiag,extracare2026},
optimal transport~\cite{ottehr2025},
inductive transfer with model fine-tuning~\cite{mutnuri2024,driever2025},
and data harmonization pipelines~\cite{purushotham2018benchmark,harutyunyan2019multitask}.
VRADA~\cite{purushotham2017vrada} introduced variational RNNs with adversarial alignment for
clinical time series but requires end-to-end training of the target model.
OTTEHR~\cite{ottehr2025} applies optimal transport to the same eICU--MIMIC databases we study,
yet performs population-level feature alignment without temporal modeling and without any
frozen-model constraint.
ExtraCare~\cite{extracare2026}, a concurrent work on interpretable clinical DA through concept-grounded
orthogonal decomposition, also operates in representation space.
In the broader time-series DA literature, established methods---CoDATS~\cite{wilson2020codats},
RAINCOAT~\cite{he2023raincoat}, CLUDA~\cite{ozyurt2023cluda},
ACON~\cite{liu2024acon}---learn domain-invariant representations through a shared feature
extractor, typically with adversarial~\cite{wilson2020codats}, contrastive~\cite{ozyurt2023cluda},
or spectral kernel~\cite{liu2021advskm} objectives.
ADATIME~\cite{ragab2023adatime} benchmarks eleven such methods across fifty scenarios, all
operating in representation space.
MAPU~\cite{ragab2023mapu} addresses the orthogonal constraint of restricted \emph{data} access
(source-free adaptation) but still modifies the target model.
A unifying limitation of every method above is the assumption that the target model's
parameters are trainable: they either fine-tune the predictor or jointly learn a feature
extractor.
When the target predictor must remain frozen---as required for regulatory compliance with
locked clinical algorithms under the FDA's Software as a Medical Device
(SaMD) framework~\cite{fda2021samd}---these methods are inapplicable,
motivating input-space adaptation approaches.

\paragraph{Frozen-model adaptation and input-space methods.}
A growing body of work modifies \emph{inputs} rather than model weights.
Prompt tuning~\cite{lester2021prompt} learns a fixed soft-token prefix for frozen language models;
Visual Prompt Tuning (VPT)~\cite{jia2022vpt} extends this to frozen vision transformers.
Adversarial reprogramming~\cite{elsayed2019reprogramming} repurposes frozen ImageNet classifiers
via a single universal additive perturbation.
Voice2Series~\cite{yang2021voice2series} reprograms frozen acoustic models for time-series
classification, and Time-LLM~\cite{jin2024timellm} reprograms frozen LLMs (GPT-2, LLaMA) for
time-series forecasting---the most prominent frozen-model work in the time-series literature,
with over 1000 citations.
TATO~\cite{qiu2026tato}, a concurrent ICLR~2026 contribution, explicitly proposes
``adapt data to model'' for frozen time-series foundation models using handcrafted
transformation operators (slicing, normalization, outlier correction).
L2C~\cite{chi2025l2c} learns input-space modifications for frozen CLIP at test time.
PixelDA~\cite{bousmalis2017pixelda} performs input-space DA in vision via GANs but does
not freeze the downstream model.
SHOT~\cite{liang2020shot} partially freezes the classifier head while adapting the feature
extractor; our constraint is stronger in that the \emph{entire} model is frozen.
Table~\ref{tab:frozen_comparison} summarizes the key distinctions.
These methods face three limitations in our setting.
First, prompt tuning, VPT, and reprogramming apply \emph{input-independent} transformations: they learn a single perturbation that is shared identically across all inputs.
While our translator also has fixed (learned) parameters, its output is \emph{input-conditioned}---each patient's time series receives a distinct transformation adapted to its specific values, length (1 to 169 timesteps), and missingness pattern across 48+ features.
A universal additive perturbation cannot provide such sample-level adaptation, as the correction needed for one patient's lab values may be entirely inappropriate for another's.
Second, the vast majority target \emph{task} adaptation (repurposing a model for a new task),
not \emph{domain} adaptation (translating data from one hospital's distribution to another's).
Third, handcrafted transforms (TATO) rely on a predefined set of operators (slicing, normalization, outlier correction) applied independently per feature.
Clinical domain shift, however, involves nonlinear, \emph{context-dependent} distributional differences: the correction required for a lab value depends on concurrent vitals, medications, and the patient's trajectory---interactions that a fixed pipeline of feature-wise operators cannot model.
To our knowledge, no prior frozen-model work simultaneously incorporates instance-level retrieval from the target domain,
evaluates on clinical EHR data, and provides theoretical grounding via gradient dynamics.
These capacity limitations motivate our
neural translator, which uses instance-level retrieval from the target domain for contextual
grounding.

\begin{table}[t]
\centering
\caption{Comparison of frozen-model and input-space adaptation methods. Our approach is the
only one that combines sample-conditional neural translation, target-domain retrieval,
variable-length support, and cross-domain DA for clinical time series.}
\label{tab:frozen_comparison}
\small
\setlength{\tabcolsep}{4pt}
\begin{tabular}{lccccc}
\toprule
\textbf{Method} & \textbf{Input modification} & \textbf{Sample-cond.} & \textbf{Retrieval} & \textbf{Var.-length} & \textbf{Purpose} \\
\midrule
Prompt Tuning~\cite{lester2021prompt} & Fixed prefix & \xmark & \xmark & \xmark & Task adapt. \\
VPT~\cite{jia2022vpt} & Prepended tokens & \xmark & \xmark & \xmark & Task adapt. \\
Voice2Series~\cite{yang2021voice2series} & Learned transform & \xmark & \xmark & \xmark & Cross-task \\
Time-LLM~\cite{jin2024timellm} & Reprogramming layer & \xmark & \xmark & Fixed patches & Cross-modality \\
TATO~\cite{qiu2026tato} & Handcrafted pipeline & Partially & \xmark & \cmark & Domain adapt. \\
\textbf{Ours} & \textbf{Neural translator} & \cmark & \cmark & \cmark & \textbf{Cross-domain DA} \\
\bottomrule
\end{tabular}
\end{table}

\paragraph{Retrieval-augmented models for time series.}
Retrieval-augmented generation (RAG)~\cite{lewis2020rag} established the paradigm of
conditioning neural models on retrieved external knowledge.
$k$NN-LM~\cite{khandelwal2020knnlm} augments a frozen language model with nearest-neighbor
retrieval from a token-level datastore, achieving domain adaptation by simply swapping the
datastore---without any model retraining.
$k$NN-MT~\cite{khandelwal2021knnmt} extends this to machine translation, performing
per-timestep $k$-NN lookup from a domain-specific datastore and interpolating
\emph{output token distributions} with the frozen model's predictions.
RETRO~\cite{borgeaud2022retro} introduced chunked cross-attention over retrieved document
passages---the closest architectural precedent for our fusion mechanism.
However, RETRO retrieves passages from the \emph{same} domain to augment language generation,
whereas we retrieve from the \emph{target} domain to transform \emph{source} inputs,
making retrieval an instrument of domain translation rather than knowledge augmentation.
In the time-series literature, retrieval augmentation is an emerging paradigm at top venues:
RATD~\cite{liu2024ratd} retrieves relevant series to guide diffusion-based forecasting
(NeurIPS~2024),
TS-RAG~\cite{ning2025tsrag} augments time-series foundation models with an adaptive
retrieval mixer (NeurIPS~2025),
and RAFT~\cite{raft2025} exploits inductive biases from retrieved examples for forecasting
(ICML~2025).
RAM-EHR~\cite{xu2024ramehr} retrieves textual medical knowledge from PubMed and
knowledge graphs to augment EHR predictions, but retrieves \emph{text}, not time-series data.
Our retrieval translator differs from all prior work in three respects:
(a)~we retrieve encoded \emph{target-domain time-series windows}---representations of what the
frozen predictor expects to see---not text passages or output distributions;
(b)~retrieved windows serve as \emph{cross-attention context} in a Transformer decoder,
enabling selective feature-level borrowing, rather than being interpolated with output
distributions ($k$NN-MT) or concatenated with the query (RAG); and
(c)~the purpose is cross-domain translation, not knowledge augmentation or forecasting.
$k$NN-MT~\cite{khandelwal2021knnmt} is the closest ancestor: both systems pair a frozen
model with per-step $k$-NN from a domain-specific store.
The critical difference is that $k$NN-MT interpolates output token distributions at
the decoder, operating within a single language, while our retrieval translator uses
retrieved embeddings as cross-attention context to transform the \emph{input}, bridging
fundamentally different data distributions (eICU vs.\ MIMIC-IV feature spaces with
heterogeneous recording practices, device calibrations, and documentation conventions).

\paragraph{Foundation models as an alternative paradigm.}
Large-scale pre-training offers a complementary path to cross-hospital generalization.
ICareFM~\cite{burger2025icarefm} pre-trains on 650K ICU stays from multiple hospitals with
a self-supervised time-to-event objective, achieving zero-shot transfer.
CLMBR/EHRSHOT~\cite{wornow2023ehrshot} provides a 141M-parameter EHR foundation model
with strong few-shot evaluation on 6,739 patients.
LLEMR~\cite{wu2024llemr} instruction-tunes an LLM on 400K+ MIMIC-IV examples.
Med-BERT~\cite{rasmy2021medbert} and BEHRT~\cite{li2020behrt} pre-train Transformers on
28.5M and 1.6M patient records, respectively.
We build directly on the YAIB benchmark framework~\cite{vandewater2024yaib}, whose
multi-center preprocessing pipeline and five standardized task definitions provide our
experimental infrastructure and frozen baselines (LSTMs in our experiments, though the input-space translation paradigm is architecture-agnostic).
Foundation models require (a)~massive multi-center datasets (100K--650K stays),
(b)~weeks of compute for pre-training, and (c)~replacing the deployed predictor
entirely, necessitating full regulatory re-validation.
Our approach is complementary, not competitive: when a hospital has a validated, deployed
predictor that cannot be retrained due to regulatory, institutional, or practical constraints,
input-space translation is the only viable adaptation strategy.
Foundation models and our translators address different points on the
\emph{deployment constraint spectrum}---from unconstrained model replacement to
fully frozen predictor preservation.

\paragraph{Domain adaptation theory and gradient dynamics.}
Ben-David et al.~\cite{bendavid2010} bound target-domain error by
$\varepsilon_T \leq \varepsilon_S + d_{\mathcal{H}\Delta\mathcal{H}}(\mathcal{S}, \mathcal{T}) + \lambda^*$,
where $d_{\mathcal{H}\Delta\mathcal{H}}$ measures distribution divergence and $\lambda^*$
is the ideal joint error.
Mansour et al.~\cite{mansour2009} refine this via discrepancy distance tailored to the
loss and hypothesis class; Zhang et al.~\cite{zhang2019mdd} bridge theory and algorithms
through margin disparity discrepancy.
Canonical methods operationalize these bounds:
DANN~\cite{ganin2016dann} minimizes $\mathcal{H}$-divergence via gradient reversal,
Deep CORAL~\cite{sun2016coral} aligns second-order feature statistics,
and DAN~\cite{long2015dan} matches multi-kernel MMD~\cite{gretton2012mmd} across layers.
All require a \emph{learnable} feature extractor.
In our frozen-model setting, representation-space alignment is impossible by construction: since $\nabla_\phi f_T = 0$, no gradient signal can reshape the frozen model's internal representations, and any alignment objective (MMD, adversarial, CORAL) applied to intermediate layers has no learnable pathway to reduce the domain gap.
We instead adapt these methods to operate in input space and evaluate them as baselines (detailed in the full paper).
%
The multi-objective structure of our training---minimizing task loss $\mathcal{L}_\text{task}$
(prediction accuracy through the frozen model) and fidelity loss $\mathcal{L}_\text{fid}$
(proximity to the source distribution) simultaneously---connects to the gradient conflict
literature in multi-task learning.
Sener and Koltun~\cite{sener2018moo} frame this as multi-objective optimization;
PCGrad~\cite{yu2020pcgrad} defines conflict as $\cos(\nabla\mathcal{L}_i, \nabla\mathcal{L}_j) < 0$
and projects conflicting gradients onto the normal plane;
CAGrad~\cite{liu2021cagrad} provides convergence guarantees under conflict.
In the DA literature specifically,
Du et al.~\cite{du2024gh} analyze obtuse-angle gradients between alignment and classification
losses in standard UDA and propose gradient rotation (GH/GH++),
while Phan et al.~\cite{phan2024pga} align per-domain gradients for prompt-based adaptation
of frozen vision-language models.
Our contribution extends gradient conflict analysis to the frozen-model DA setting.
We show that the cosine similarity
$\cos(\nabla_\theta\mathcal{L}_\text{task}, \nabla_\theta\mathcal{L}_\text{fid})$
between task and fidelity gradients---measured at the translator parameters $\theta$---is
a diagnostic that predicts method--task suitability:
$+0.84$ for mortality (aligned; delta translation suffices) versus $-0.21$ for sepsis
(conflicting; delta translation fails catastrophically).
Crucially, where PCGrad and CAGrad intervene \emph{algorithmically} by manipulating gradient
directions, our retrieval translator resolves the conflict \emph{architecturally}: the
cross-attention pathway over target-domain windows provides an alternative information
channel that reduces the dependence of the task gradient on input-space fidelity,
substantially mitigating the destructive interference without requiring gradient surgery
(detailed empirical gradient measurements in the full paper).

\medskip
\noindent
To summarize, no prior work combines a frozen target predictor, learned neural input-space
translation, instance-level target-domain retrieval, and gradient alignment theory for
clinical time-series domain adaptation.
The frozen-model constraint is motivated by real-world regulatory requirements, yet it
renders all existing clinical DA methods, general time-series DA methods, and foundation
model approaches inapplicable without fundamental modification.
Our retrieval-augmented translator fills this gap: it operates entirely in input space,
grounds translations in target-domain examples, and provides a theoretical framework
that explains when and why the approach succeeds.
